\section{Dual Front-End Architecture}
\label{sec:frontend}

The compiler exposes two input paths that converge on the same internal pipeline.
This dual front-end is the architectural key to serving both human authors
(who prefer text DSLs) and AI agents (who prefer structured data).

\subsection{The Two Input Paths}

\ourdiagram[0.95\linewidth]{dual-frontend}{%
  Dual front-end: Mermaid DSL text (Path~A) passes through a PEG/Nearley parser;
  structured IR JSON/YAML (Path~B) passes through schema validation.
  Both paths converge at the same Domain IR, then proceed through the shared
  layout and rendering pipeline to SVG, PNG, and PDF outputs.}

\paragraph{Path A: Mermaid-Superset DSL.}
Human authors write diagrams in Mermaid syntax (or our extensions).
A parser converts the text to the grammar's Domain IR.
This path preserves the ``diagrams-as-code'' ergonomics that made Mermaid popular.

\paragraph{Path B: Structured IR.}
Agents emit JSON or YAML conforming to the Domain IR's JSON Schema.
No parsing step is needed---only schema validation.
This path enables grammar-constrained generation, round-trip editing, and tool integration.

\subsection{Convergence at Domain IR}

Both paths produce the same Domain IR. From that point forward, the pipeline is identical:
Domain IR $\to$ Layout Engine $\to$ Scene IR $\to$ Backend $\to$ Output.

This convergence guarantees:
\begin{itemize}
  \item A Mermaid file and its equivalent YAML produce \textbf{byte-identical output}.
  \item Agents can consume human-authored diagrams (parse DSL $\to$ IR $\to$ emit structured).
  \item Humans can inspect agent-generated diagrams (structured IR $\to$ pretty-print DSL).
  \item Tooling (validation, linting, diffing) works on the common Domain IR regardless of input path.
\end{itemize}

\subsection{The Agent IR as First-Class API}

The structured IR is not a secondary or internal representation---it is a \textbf{first-class
public API}. Design commitments:

\begin{description}
  \item[Published JSON Schema] Every grammar's Domain IR has a versioned JSON Schema,
        suitable for OpenAI Structured Outputs, XGrammar~\cite{dong2024xgrammar},
        or llama.cpp GBNF constrained generation.
  \item[Stable contract] The IR schema follows semver. Breaking changes require a major version bump.
  \item[Round-trip fidelity] DSL $\to$ IR $\to$ DSL preserves semantics (formatting may differ).
  \item[Diff-friendly] Changes to one entity do not cascade to unrelated IR fields.
  \item[Small schemas] Each grammar's schema fits comfortably in an LLM context window
        (typically under 200 lines of JSON Schema).
\end{description}

\subsection{Parser Architecture}

The DSL parser is grammar-specific. Each grammar defines:
\begin{itemize}
  \item A PEG or Nearley grammar specification for its Mermaid-compatible syntax
  \item A CST-to-Domain-IR lowering pass
  \item Extension hooks for superset syntax (new keywords, extended directives)
\end{itemize}

The parser layer is the \textbf{new architectural component} introduced by the Mermaid-superset
positioning. It sits above the existing Domain IR / Layout Engine / Scene IR pipeline and does
not affect the proven kernel or grammar internals.

\subsection{Implementation Status}

\textbf{Built:} The structured IR path (Path B) is fully implemented for all four shipped
grammars (timeline, flow, sequence, tree). Schema validation, layout, and rendering work
end-to-end.

\textbf{Planned:} The Mermaid-superset DSL parsers (Path A) are the primary Tier-0 deliverable.
Parsers for flowchart, sequence, gantt/timeline, and mindmap syntax will wire the existing
Domain IRs to Mermaid text input.
